\documentclass[11pt]{article}

\usepackage{amsmath}
\usepackage[margin=1in]{geometry}
\usepackage{booktabs}
\usepackage{array}
\usepackage{hyperref}
\usepackage{setspace}
\usepackage{xcolor}
\usepackage{caption}

\title{Supplementary Materials\\
\large Calibration, Uncertainty Communication, and Deployment Readiness\\
in CKD Risk Prediction: A Framework Evaluation Study}

\author{Michael Eniolade\\
        \small University of the Cumberlands\\
        \small \href{mailto:meniolade20593@ucumberlands.edu}{meniolade20593@ucumberlands.edu}}

\date{}

\begin{document}
\maketitle

\tableofcontents
\clearpage

% ─────────────────────────────────────────────────────────────────────────────
\section*{S1 — Hyperparameter Search Grids}
\addcontentsline{toc}{section}{S1 — Hyperparameter Search Grids}

All classifiers were tuned using 5-fold stratified cross-validation on the
UCI training fold ($n = 279$) with AUROC as the scoring metric.
\texttt{GridSearchCV} (scikit-learn 1.5.2) was used throughout.
Random seed: 42.

\subsection*{Logistic Regression (LR)}

Solver: \texttt{lbfgs}; max iterations: 1000; class weight: balanced.

\begin{center}
\begin{tabular}{ll}
\toprule
Hyperparameter & Values searched \\
\midrule
$C$ (inverse regularization strength) & 0.001, 0.01, 0.1, 1, 10, 100 \\
\bottomrule
\end{tabular}
\end{center}

\noindent\textbf{Selected:} $C = 0.001$\quad\textbf{CV AUROC:} 1.0000

\subsection*{Random Forest (RF)}

\begin{center}
\begin{tabular}{ll}
\toprule
Hyperparameter & Values searched \\
\midrule
\texttt{n\_estimators}    & 100, 200 \\
\texttt{max\_depth}       & 5, 10, 20, None \\
\texttt{min\_samples\_split} & 2, 5, 10 \\
\texttt{max\_features}    & ``sqrt'', ``log2'' \\
\bottomrule
\end{tabular}
\end{center}

\noindent\textbf{Selected:} \texttt{n\_estimators}=100, \texttt{max\_depth}=20,
\texttt{min\_samples\_split}=10, \texttt{max\_features}=``sqrt''\quad
\textbf{CV AUROC:} 0.9997

\subsection*{XGBoost (XGB)}

\begin{center}
\begin{tabular}{ll}
\toprule
Hyperparameter & Values searched \\
\midrule
\texttt{n\_estimators}      & 100, 200 \\
\texttt{max\_depth}         & 3, 5, 7 \\
\texttt{learning\_rate}     & 0.05, 0.1, 0.2 \\
\texttt{subsample}          & 0.8, 1.0 \\
\texttt{colsample\_bytree}  & 0.6, 0.8, 1.0 \\
\texttt{gamma}              & 0, 0.1, 0.5 \\
\bottomrule
\end{tabular}
\end{center}

\noindent\textbf{Selected:} \texttt{n\_estimators}=100, \texttt{max\_depth}=3,
\texttt{learning\_rate}=0.2, \texttt{subsample}=0.8,
\texttt{colsample\_bytree}=0.6, \texttt{gamma}=0.5\quad
\textbf{CV AUROC:} 0.9992

\subsection*{Support Vector Machine (SVM)}

\texttt{probability=True} enabled internally. Class weight: balanced.

\begin{center}
\begin{tabular}{ll}
\toprule
Hyperparameter & Values searched \\
\midrule
$C$            & 0.1, 1, 10, 100 \\
\texttt{kernel}& ``rbf'', ``poly'' \\
\texttt{gamma} & ``scale'', 0.1, 1.0 \\
\bottomrule
\end{tabular}
\end{center}

\noindent\textbf{Selected:} $C=1$, \texttt{kernel}=``poly'',
\texttt{gamma}=0.1\quad\textbf{CV AUROC:} 1.0000

\subsection*{Gaussian Naive Bayes (NB)}

\begin{center}
\begin{tabular}{ll}
\toprule
Hyperparameter & Values searched \\
\midrule
\texttt{var\_smoothing} & $10^{-11}$, $10^{-9}$, $10^{-7}$, $10^{-5}$, $10^{-3}$ \\
\bottomrule
\end{tabular}
\end{center}

\noindent\textbf{Selected:} \texttt{var\_smoothing}=$10^{-11}$\quad
\textbf{CV AUROC:} 1.0000

\subsection*{Post-hoc Calibration Methods}

Applied after model selection using
\texttt{sklearn.calibration.CalibratedClassifierCV} with
\texttt{FrozenEstimator} to prevent base-model refitting.

\begin{center}
\begin{tabular}{ll}
\toprule
Method & Parameters \\
\midrule
Platt scaling (sigmoid) & \texttt{method="sigmoid"}, \texttt{cv="prefit"} \\
Isotonic regression     & \texttt{method="isotonic"}, \texttt{cv="prefit"} \\
\bottomrule
\end{tabular}
\end{center}

Calibration fitting used the UCI validation fold ($n=60$, indices in
\texttt{data/processed/uci\_splits.json}). Best variant per model selected
by lowest ECE on the UCI test fold ($n=61$).

\clearpage

% ─────────────────────────────────────────────────────────────────────────────
\section*{S2 — Reliability Diagrams}
\addcontentsline{toc}{section}{S2 — Reliability Diagrams}

Reliability diagrams (calibration curves) are provided for all
model--dataset combinations. Each diagram plots mean predicted probability
($x$-axis) against fraction of positive cases per bin ($y$-axis); a
perfectly calibrated model follows the diagonal. Bins: 10 equal-width
bins. Figures generated at 300 DPI.

Figures \ref{fig:s2a} and \ref{fig:s2b} in the main manuscript display the
combined pre- and post-calibration reliability diagrams on the UCI test set.
Figure \ref{fig:s2c} in the main manuscript displays MIMIC-IV reliability
diagrams for the best-calibrated variant of each model.

\subsection*{UCI Test Set — All Variants}

\begin{center}
\begin{tabular}{lll}
\toprule
File & Model & Variant \\
\midrule
\texttt{reliability\_compare\_LR.png}  & Logistic Regression & Base, Platt, Isotonic \\
\texttt{reliability\_compare\_RF.png}  & Random Forest       & Base, Platt, Isotonic \\
\texttt{reliability\_compare\_XGB.png} & XGBoost             & Base, Platt, Isotonic \\
\texttt{reliability\_compare\_SVM.png} & SVM                 & Base, Platt, Isotonic \\
\texttt{reliability\_compare\_NB.png}  & Naive Bayes         & Base, Platt, Isotonic \\
\bottomrule
\end{tabular}
\end{center}

\subsection*{MIMIC-IV Stress-Test Cohort — Best-Calibrated Variant}

\begin{center}
\begin{tabular}{llll}
\toprule
File & Model & Best variant & MIMIC ECE \\
\midrule
\texttt{reliability\_mimic\_LR.png}  & Logistic Regression & Isotonic & 0.761 \\
\texttt{reliability\_mimic\_RF.png}  & Random Forest       & Isotonic & 0.753 \\
\texttt{reliability\_mimic\_XGB.png} & XGBoost             & Isotonic & 0.680 \\
\texttt{reliability\_mimic\_SVM.png} & SVM                 & Isotonic & 0.755 \\
\texttt{reliability\_mimic\_NB.png}  & Naive Bayes         & Base     & 0.753 \\
\bottomrule
\end{tabular}
\end{center}

All MIMIC reliability diagrams show severe deviation above the diagonal,
consistent with models trained on a 62.5\% CKD prevalence population being
applied to a 23.7\% prevalence cohort. Each figure includes an inset
reporting ECE, MCE, Brier Score, and AUROC.

\clearpage

% ─────────────────────────────────────────────────────────────────────────────
\section*{S3 — Subgroup Calibration on MIMIC-IV}
\addcontentsline{toc}{section}{S3 — Subgroup Calibration on MIMIC-IV}

Subgroup-level ECE was computed on the MIMIC-IV demo stress-test cohort
($n = 97$) across three stratification variables: age group, diabetes
mellitus (DM) status, and hypertension (HTN) status. ECE used 5
equal-width bins to ensure stable estimates in small subgroups. Groups
with fewer than 10 patients are marked ``---'' and excluded from
criterion scoring.

\subsection*{Subgroup ECE Values}

\begin{center}
\begin{tabular}{lccccc}
\toprule
Model & Age $<65$ ($n=55$) & Age $\geq 65$ ($n=42$) & DM=Yes & HTN=Yes & Overall ($n=97$) \\
\midrule
LR  & 0.852 & 0.643 & --- & --- & 0.761 \\
RF  & 0.836 & 0.643 & --- & --- & 0.753 \\
XGB & 0.754 & 0.605 & --- & --- & 0.680 \\
SVM & 0.844 & 0.639 & --- & --- & 0.755 \\
NB  & 0.818 & 0.667 & --- & --- & 0.753 \\
\bottomrule
\end{tabular}
\end{center}

DM=Yes and HTN=Yes subgroups each fell below the 10-patient minimum
threshold after cohort construction and are excluded from ECE computation.

\subsection*{Subgroup ECE Gap (Criterion 7)}

The deployment readiness criterion requires maximum ECE gap across valid
subgroups $\leq 0.05$. The gap is max(subgroup ECE) $-$ min(subgroup ECE).

\begin{center}
\begin{tabular}{lcccc}
\toprule
Model & Max ECE & Min ECE & Gap & Criterion 7 \\
\midrule
LR  & 0.852 & 0.643 & 0.209 & FAIL \\
RF  & 0.836 & 0.643 & 0.194 & FAIL \\
XGB & 0.754 & 0.605 & 0.148 & FAIL \\
SVM & 0.844 & 0.639 & 0.205 & FAIL \\
NB  & 0.818 & 0.605 & 0.213 & FAIL \\
\bottomrule
\end{tabular}
\end{center}

All models fail the subgroup equity criterion. Younger patients (age $<65$)
show systematically worse calibration than older patients across all five
models, ECE gaps 0.15--0.21. This pattern is a secondary effect of the
overall calibration collapse: models overestimate CKD probability for all
patients, and overestimation is more pronounced in younger patients who have
lower true CKD prevalence in the demo cohort. Even the best subgroup
(age $\geq 65$, ECE $\approx 0.64$) remains far above the acceptable
threshold of ECE $\leq 0.10$.

\clearpage

% ─────────────────────────────────────────────────────────────────────────────
\section*{S4 — Reproducibility Instructions}
\addcontentsline{toc}{section}{S4 — Reproducibility Instructions}

All results are fully reproducible from raw data. The pipeline runs in
approximately 5--10 minutes on a standard laptop.

\subsection*{Software Environment}

\begin{center}
\begin{tabular}{ll}
\toprule
Package & Version \\
\midrule
Python           & 3.11+ \\
scikit-learn     & 1.5.2 \\
xgboost          & 2.1.1 \\
imbalanced-learn & 0.12.4 \\
netcal           & 1.3.5 \\
mapie            & 0.8.6 \\
pandas           & 2.2+ \\
numpy            & 1.26+ \\
matplotlib       & 3.9+ \\
joblib           & 1.4+ \\
\bottomrule
\end{tabular}
\end{center}

Install all dependencies: \quad\texttt{pip install -r requirements.txt}

\subsection*{Data Acquisition}

\textbf{UCI CKD dataset.} Download \texttt{chronic\_kidney\_disease\_full.arff}
from the UCI Machine Learning Repository (no account required):
\url{https://archive.ics.uci.edu/dataset/336/chronic+kidney+disease}.
Place the file in \texttt{first\_paper/data/raw/}.

\textbf{MIMIC-IV Clinical Database Demo v2.2.} Download from PhysioNet
(free account required; no credentialing needed for the demo release):
\url{https://physionet.org/content/mimic-iv-demo/2.2/}.
Extract to any local directory and set the path in
\texttt{code/01\_data\_prep/extract\_mimic\_local.py}.

\subsection*{Reproduction Steps}

Run each script from the \texttt{first\_paper/} directory in the
following order. All paths are relative to that directory.

\begin{enumerate}
\item \texttt{python code/01\_data\_prep/clean\_uci.py}\\
      Output: \texttt{data/processed/uci\_ckd\_clean.csv},
      \texttt{data/processed/uci\_splits.json}

\item \texttt{python code/01\_data\_prep/extract\_mimic\_local.py}\\
      Output: \texttt{data/processed/mimic\_ckd\_cohort.csv}

\item \texttt{python code/01\_data\_prep/harmonize\_mimic.py}\\
      Output: \texttt{data/processed/mimic\_ckd\_clean.csv},
      \texttt{data/processed/schema\_comparison.csv}

\item \texttt{python code/02\_modeling/train\_models.py}\\
      Output: saved model \texttt{.pkl} files; Tables T1 and metadata CSVs

\item \texttt{python code/03\_calibration/calibrate\_models.py}\\
      Output: calibrated model files; pre-/post-calibration reliability
      diagrams; Table T2 partial

\item \texttt{python code/03\_calibration/external\_calibration.py}\\
      Output: MIMIC reliability diagrams; Table T2 complete

\item \texttt{python code/03\_calibration/bootstrap\_ci.py}\\
      Output: 95\% bootstrap CI columns merged into Table T2

\item \texttt{python code/04\_uncertainty/conformal\_uci.py}\\
      Output: conformal model files; Table T3 partial; Figures F3a/F3b

\item \texttt{python code/04\_uncertainty/mimic\_conformal.py}\\
      Output: Table T3 complete

\item \texttt{python code/05\_deployment\_checklist/subgroup\_calibration.py}\\
      Output: subgroup ECE table (S3)

\item \texttt{python code/05\_deployment\_checklist/score\_checklist.py}\\
      Output: Table T4; Figure F4
\end{enumerate}

\subsection*{Expected Key Metrics}

\textbf{Table T2 — Calibration summary (ECE and AUROC):}

\begin{center}
\begin{tabular}{lcccc}
\toprule
Model & UCI ECE & UCI AUROC & MIMIC ECE & MIMIC AUROC \\
\midrule
LR  & 0.022 & 1.000 & 0.761 & 0.485 \\
RF  & 0.000 & 1.000 & 0.753 & 0.507 \\
XGB & 0.007 & 1.000 & 0.680 & 0.579 \\
SVM & 0.015 & 1.000 & 0.755 & 0.483 \\
NB  & 0.016 & 1.000 & 0.753 & 0.477 \\
\bottomrule
\end{tabular}
\end{center}

\textbf{Table T3 — Conformal prediction coverage:}

\begin{center}
\begin{tabular}{lcc}
\toprule
Model & UCI Coverage & MIMIC Coverage \\
\midrule
LR  & 0.803 & 0.237 \\
RF  & 0.967 & 0.206 \\
XGB & 0.967 & 0.227 \\
SVM & 0.918 & 0.216 \\
NB  & 0.984 & 0.247 \\
\bottomrule
\end{tabular}
\end{center}

\subsection*{Random Seeds}

All stochastic operations use \texttt{random\_state=42}: train/test
splitting, cross-validation folds, Random Forest, XGBoost, and bootstrap
resampling. Results are bit-for-bit identical across runs on the same
machine and package versions.

\end{document}
